\section{半单李代数的极大子代数}

定理 1。令 V 为有限维向量空间，g 为 \(\mathfrak{gl}(V)\) 中的约化李子代数，q 为 g 的李子代数，\(\Phi\) 为 \((x,y)\mapsto\mathrm{Tr}(xy)\) 上的双线性型 \(g\times g\)。假设 q 关于 \(\Phi\) 的正交补 n 是 g 的李子代数，且由 V 的幂零自同态组成。则 q 是 g 的抛物子代数。

\begin{enumerate}
\def\labelenumi{\alph{enumi})}
\tightlist
\item
  \(\mathfrak{q}\) 是 \(\mathfrak{n}\) 在 \(\mathfrak{g}\) 中的正规化子：令 \(\mathfrak{p}\) 为这个正规化子。令 \(x \in \mathfrak{q}\) 且 \(y \in \mathfrak{n}\)；对所有 \(z \in \mathfrak{q}\)，都有 \([z, x] \in \mathfrak{q}\)，所以
\end{enumerate}

\[
\Phi ([ x, y ], z) = \Phi (y, [ z, x ]) = 0;
\]

换言之，\([x,y] \in n\)。因此 \(q \subset p\)。由于 n 是 p 的理想且由 V 的幂零自同态组成，P 关于 \(\varPhi\) 与 n 正交（第 I 章第3号命题4 d））。由于 \(\varPhi\) 非退化 \(^{4}\)，\(p \subset q\)，故断言得证。

\begin{enumerate}
\def\labelenumi{\alph{enumi})}
\setcounter{enumi}{1}
\tightlist
\item
  存在 \(\mathfrak{gl}(V)\) 中的约化李子代数 m，使得 q 是 m 与 n 的半直积：令 \(\mathfrak{n}_{\mathrm{V}}(\mathfrak{q})\) 为 q 中由 V 的幂零自同态组成的最大理想。则 \(\mathfrak{n}_{\mathrm{V}}(\mathfrak{q})\) 包含 n，且它与 q 正交（同上）；因此 \(\mathfrak{n} = \mathfrak{n}_{\mathrm{V}}(\mathfrak{q})\)。此外，按假设，g 是 \(\mathfrak{gl}(V)\) 中的约化李代数，因而是可分解的（第 VII 章第5节第1号命题2）；由于 q 是 g 与 \(\mathfrak{gl}(V)\) 中 n 的正规化子的交，它是一个可分解李代数（同上，命题3的推论1）。因此断言由第 VII 章第5节第3号命题7 得到。
\end{enumerate}

取 \(\mathfrak{h}\) 的一个嘉当子代数 \(\mathfrak{m}\)；记 \(\mathfrak{g}_1\) 为 \(\mathfrak{h}\) 在 \(\mathfrak{g}\) 中的交换子，并置 \(\mathfrak{q}_1 = \mathfrak{q} \cap \mathfrak{g}_1, \mathfrak{n}_1 = \mathfrak{n} \cap \mathfrak{g}_1\)。

\begin{enumerate}
\def\labelenumi{\alph{enumi})}
\setcounter{enumi}{2}
\tightlist
\item
  李代数 \(\mathfrak{g}_1, \mathfrak{q}_1\) 和 \(\mathfrak{n}_1\) 满足与 \(\mathfrak{g}, \mathfrak{q}\) 和 \(\mathfrak{n}\) 相同的假设：由于 \(\mathfrak{m}\) 是 \(\mathfrak{gl}(\mathrm{V})\) 中的约化李代数，\(\mathfrak{h}\) 是交换的，并由 V 的半单自同态组成（第 VII 章第2节第4号定理2的推论3）。因此 \(\mathfrak{g}_1 = \mathfrak{g}^0(\mathfrak{h})\) 是 \(\mathfrak{g}\) 中的约化李代数（第 VII 章第1节第3号命题11），从而也是 \(\mathfrak{gl}(\mathrm{V})\) 中的约化李代数（第 I 章第6节第6号命题7的推论2）。显然，\(\mathfrak{n}_1\) 由 V 的幂零自同态组成。由于 \(\mathfrak{h}\) 是 \(\mathfrak{q}\) 的子代数，且在 \(\mathfrak{gl}(V)\) 中约化，\(\mathfrak{h}\) 在 \(\mathfrak{q}\) 上的伴随表示是半单的；按构造，\(\mathfrak{q}_1\) 是 \(\mathrm{ad}_{\mathfrak{q}}(\mathfrak{h})\) 的不变量集合，所以 \(\mathfrak{q} = \mathfrak{q}_1 + [\mathfrak{h}, \mathfrak{q}]\)（第 I 章第3节第5号命题6）。由于
\end{enumerate}

\[
\Phi (\mathfrak {g} _ {1}, [ \mathfrak {h}, \mathfrak {q} ]) = \Phi ([ \mathfrak {h}, \mathfrak {g} _ {1} ], \mathfrak {q}) = 0,
\]

\(\mathfrak{g}_1\) 中的元素与 \(\mathfrak{q}_1\) 正交，当且仅当它与 \(\mathfrak{q}\) 正交；因此，\(\mathfrak{n}_1 = \mathfrak{g}_1 \cap \mathfrak{n}\) 是 \(\mathfrak{q}_1\) 中 \(\mathfrak{g}_1\) 的正交补。

\begin{enumerate}
\def\labelenumi{\alph{enumi})}
\setcounter{enumi}{3}
\item
   的嘉当子代数 \(\mathfrak{h}\) 是 \(\mathfrak{m}\)\(\mathfrak{g}\) 的嘉当子代数：我们有 \(\mathfrak{q} = \mathfrak{m} \oplus \mathfrak{n}\) 且 \(\mathfrak{h} = \mathfrak{m} \cap \mathfrak{g}_1\)，所以立即得到 \(\mathfrak{q}_1 = \mathfrak{h} \oplus \mathfrak{n}_1\)。此外，\([\mathfrak{h}, \mathfrak{n}_1] = 0\)，\(\mathfrak{h}\) 交换且 \(\mathfrak{n}_1\) 幂零，所以李代数 \(\mathfrak{q}_1\) 幂零。由 a）和 c），\(a)\)\(c)\)\(\mathfrak{q}_1\) 是 \(\mathfrak{n}_1\) 在 \(\mathfrak{g}_1\) 中的正规化子；更不用说，\(\mathfrak{q}_1\) 等于它在 \(\mathfrak{g}_1\) 中的正规化子，因而是 \(\mathfrak{g}_1\) 的嘉当子代数。由于 \(\mathfrak{g}_1\) 是 \(\mathfrak{gl}(V)\) 中的约化李代数，由第 VII 章第2节第4号定理2的推论3 可知，\(\mathfrak{q}_1\) 由 V 的半单自同态组成；于是，由于 \(\mathfrak{n}_1\) 由 V 的幂零自同态组成，有 \(\mathfrak{n}_1 = 0\)。因此，\(\mathfrak{h} = \mathfrak{q}_1\) 是 \(\mathfrak{g}_1\) 的嘉当子代数，并且由于 \(\mathfrak{g}_1\) 规范化 \(\mathfrak{h}\)，有 \(\mathfrak{h} = \mathfrak{g}_1\)。于是我们证明了 \(\mathfrak{h}\) 的每个元素都是 \(\mathfrak{g}\) 的半单元素，且 \(\mathfrak{h}\) 在 \(\mathfrak{g}\) 中的交换子等于 \(\mathfrak{h}\)；由此 \(\mathfrak{h} = \mathfrak{g}^0(\mathfrak{h})\)，所以 \(\mathfrak{h}\) 是 \(\mathfrak{g}\) 的嘉当子代数。
\item
  \(\mathfrak{q}\) 是 \(\mathfrak{g}\) 的抛物子代数：由上文，\(\mathfrak{h}\) 是 \(\mathfrak{g}\) 的嘉当子代数，\(\mathfrak{n}\) 由 \(\mathfrak{g}\) 的幂零元素组成，且 \([\mathfrak{h},\mathfrak{n}]\subset \mathfrak{n}\)。令 \(\bar{k}\) 为 k 的代数闭包；按定义，q 在 g 中是抛物的，当且仅当 \(k\)\(\mathfrak{q}\)\(\mathfrak{g}\)\(\bar{k}\otimes_k\mathfrak{q}\) 是 \(\bar{k}\otimes_k\mathfrak{g}\) 的抛物子代数。上述性质在标量扩张下保持不变，因此证明时可以限于 \(\mathfrak{h}\) 分裂的情形。令 R 为 \((\mathfrak{g},\mathfrak{h})\) 的根系；由第3节第1号命题2 (v)，存在 R 的子集 P，使得 \(P\cap (-P) = \varnothing\) 且 \(\mathfrak{n} = \sum_{\alpha \in P}\mathfrak{g}^{\alpha}\)。令 \(P'\) 为满足  的根 \(\alpha\) 所成的集合；有 \(-\alpha\notin P\)\(P' \cup (-P') = R\)，而  中 n 的正交补 \(\mathfrak{q}\) 等于 \(\mathfrak{n}\)\(\mathfrak{g}\)\(\mathfrak{h} + \sum_{\alpha \in P'}\mathfrak{g}^{\alpha}\)。我们已经证明 \(\mathfrak{q}\) 是抛物的。证毕。
\end{enumerate}

引理 1。令 g 为半单李代数，V 为有限维向量空间，\(\rho\) 为 g 在 V 上的线性表示，D 为 V 的向量子空间，h 为 g 的嘉当子代数，s（分别为 \(s'\)）为满足 \(x \in h\)（分别为 \(\rho(x)\mathrm{D} \subset \mathrm{D}\)）的 \(\rho(x)\mathrm{D} = 0\) 所成的集合，\(\Phi\) 为与表示 \(^{5}\) 相关联的 \(\rho\) g 上的双线性型。

\begin{enumerate}
\def\labelenumi{(\roman{enumi})}
\item
  若 \(\mathfrak{h}\) 分裂，则 \(\mathfrak{s}\) 和 \(\mathfrak{s}'\) 是 \(\mathfrak{h}\) 中关于 \(\mathbf{Q}\) 有理的向量子空间。
\item
  若 \(\rho\) 是单射，则 \(\Phi\) 在 \(\mathfrak{s}\)（分别为 \(\mathfrak{s}'\)）上的限制是非退化的。
\end{enumerate}

假设嘉当子代数 \(\mathfrak{h}\) 是分裂的。令 \(d\) 为 D 的维数；置 \(\mathrm{W} = \bigwedge^{d}(\mathrm{V})\) 和 \(\sigma = \bigwedge^{d}(\rho)\)；还以 \((e_1, \ldots, e_d)\) 表示 D 的一个基，并令 \(e = e_1 \wedge \cdots \wedge e_d\) 为与 D 相关的可分解 \(d\)-向量。令 P 为 \(\sigma\) 关于 \(\mathfrak{h}\) 的权集合；记 \(\mathrm{W}^{\mu}\) 为

与权 \(\mu\) 对应的 W 的子空间，并置 \(e = \sum_{\mu \in P} e^{\mu}\)（其中对所有 ，均有 \(e^{\mu} \in W^{\mu}\)）；最后，令 P' 为满足 \(\mu \in P\)\(P'\) 的权 \(\mu\) 所成的集合，令 P'' 为 P' 中元素的差集。令 x 属于 h；则 x 属于 s，当且仅当存在 c ∈ k 使 \(e^{\mu} \neq 0\)\(P''\)\(P'\)\(\rho(x).e = c.e\)（第 VII 章第5节第4号引理2 (i)）。由于 \(\rho(x).e^{\mu} = \mu(x).e^{\mu}\)，可见 \(x \in s\) 等价于关系“对所有 ，\(\mu(x) = 0\)”。现在，h 的 Q-结构是由余根 \(\mu \in P''\) 生成的 h 的 Q-向量子空间 \(h_{Q}\)，且 P'' 中的所有 \(H_{\alpha}\)\(\mu\) 在 \(P''\)\(h_{Q}\) 上取有理值；由此（《代数》第 II 章第8节第4号命题5），s 是 h 中关于 Q 有理的子空间。

对任意权 \(\mu\in P\)，令 \(p_{\mu}\) 为与分解  相对应的到 \(V^{\mu}\) 的投影；记 \(V=\bigoplus_{\mu\in P}V^{\mu}\)\(P_{1}\) 为满足  的 \(\mu\in P\) 所成的集合。显然，s' 是 h 中 P＿\(p_{\mu}(D)\neq0\)\(s'\)\(P_{1}\)\(s'\)1 各元素的核的交；如同 s 的情形一样可知，s' 是 h 中关于 Q 有理的子空间。这证明了 (i)。

通过标量扩张，只需在 k 为代数闭域、因而 \(k\)\(\mathfrak{h}\) 分裂时证明 (ii)。令 \(\mathfrak{m}\) 为 h 中关于 \(\mathfrak{h}\)\(\mathbf{Q}\) 有理的向量子空间；对 \(x\)\(\mathfrak{m}_{\mathbf{Q}} = \mathfrak{m} \cap \mathfrak{h}_{\mathbf{Q}}\) 中的每个非零 x，由 §7 第1号命题1的推论，有 \(\varPhi(x, x) > 0\)。\(\varPhi\) 在 \(\mathfrak{m}_{\mathbf{Q}}\) 上的限制非退化，从而由于  与  典则同构，\(\varPhi\) 在 \(\mathfrak{m}\) 上的限制也非退化。\(\mathfrak{m}\)\(k \otimes_{\mathbf{Q}} \mathfrak{m}_{\mathbf{Q}}\)

定义 1。令 q 为半单李代数 g 的一个李子代数。若对每个 \(x \in q\)，x 在 g 中的半单分量和幂零分量都属于 q，则称 q 在 g 中可分解。记 \(\mathfrak{n}_{\mathfrak{g}}(\mathfrak{q})\) 为 q 的根基中满足 \(ad_{g}x\) 幂零的元素 x 所成的集合。

令 \(\rho\) 为 g 在有限维向量空间 V 上的一个单射表示。我们知道（第 I 章第6节第3号定理3），g 的元素 x 是半单的（分别为幂零的），当且仅当 V 的自同态 \(\mathfrak{g}\)\(x\)\(\mathfrak{g}\)\(\rho(x)\) 是半单的（分别为幂零的）。由此立即得到，代数 \(\mathfrak{q}\) 在 \(\mathfrak{g}\) 中可分解，当且仅当 \(\rho(\mathfrak{q})\) 是 \(\mathfrak{gl}(V)\) 中按第 VII 章第5节第1号定义1 所称的可分解子代数。沿用第 VII 章第5节第3号的记号，还有

\[
\rho (\mathfrak {n} _ {\mathfrak {g}} (\mathfrak {q})) = \mathfrak {n} _ {\mathrm{V}} (\rho (\mathfrak {q})).
\]

定理 2。令 g 为半单李代数，n 为由幂零元素组成的 g 的子代数，q 为 n 在 g 中的正规化子。假设 n 正好是 q 的根基中的幂零元素所成的集合。则 q 是抛物的。

首先注意，q 是可分解的（第 VII 章第5节第1号命题3的推论1）。由定理1，只需证明 q 是关于 g 的 Killing 型 \(n^{0}\) 的 n 的正交补 \(\varPhi\)。我们知道 \(q \subset n^{0}\)（第 I 章第4节第3号命题4 d））。由第 VII 章第5节第3号命题7，存在 q 的一个在 g 中约化的子代数 m，使得 q 是 m 与 n 的半直积。~我们证明 \(\varPhi\) 在 m 上的限制非退化。令 c 为 m 的中心。由第 I 章第5节第5号命题5，有 \(\Phi([\mathfrak{m},\mathfrak{m}],\mathfrak{c})=0\)；而由第 I 章第6节第1号命题1，\(\Phi\) 在 \([m,m]\) 上的限制非退化。还需证明 \(\Phi\) 在 c 上的限制非退化。令 k 为 \([m,m]\) 的一个嘉当子代数；于是 \(k\oplus c\) 是交换的且在 g 中约化。令 h 为包含 \(k\oplus c\) 的 g 的一个嘉当子代数（第 VII 章第2节第3号命题10）。于是 \(h\cap q\) 是 q 的一个交换子代数，包含 \(k\oplus c\)，并且对每个 \(ad_{q}x\)，\(x\in h\cap q\) 是半单的；因此 \(h\cap q\) 包含于 q 的一个嘉当子代数 \(h'\) 中（第 VII 章第2节第3号命题10）；令 f 为以 n 为核的 q 到 m 的投影；则 \(f(h')\) 是 m 的一个嘉当子代数（第 VII 章第2节第1号命题4的推论2），包含 \(k\oplus c\)，所以等于 \(k\oplus c\)；这证明了 \(f(h\cap q)=k\oplus c\)，并且由于 h 的每个元素在 g 中都是半单的，有 \(h\cap q=k\oplus c\)。于是，

\[
\mathfrak {c} = \{x \in \mathfrak {h} \mid [ x, \mathfrak {n} ] \subset \mathfrak {n} \text {   and   } [ x, [ \mathfrak {m}, \mathfrak {m} ] ] = 0 \}.
\]

由引理1，\(\varPhi\) 在 \(\mathfrak{c}\) 上的限制非退化。

令 \(\mathfrak{q}^0\) 为 \(\mathfrak{q}\) 中相对于 \(\mathfrak{g}\) 的 \(\varPhi\) 的正交补。上文证明了 \(\mathfrak{q} \cap \mathfrak{q}^0 = \mathfrak{n}\)。假设 \(\mathfrak{q} \neq \mathfrak{q}^0\)，于是 \(\mathfrak{q}^0 \neq \mathfrak{n}\)（且 \(\mathfrak{q}^0 \supset \mathfrak{n}\)）。由于 \(\operatorname{ad}_{\mathfrak{g}} \mathfrak{n}\) 保持 \(\mathfrak{q}\) 不变，\(\operatorname{ad}_{\mathfrak{g}} \mathfrak{n}\) 也保持 \(\mathfrak{q}^0\) 不变；Engel 定理说明存在 \(x \in \mathfrak{q}^0\)，使得 \(x \notin \mathfrak{n}\) 且 \([x, \mathfrak{n}] \subset \mathfrak{n}\)。但此时 \(x \in \mathfrak{q}^0 \cap \mathfrak{q} = \mathfrak{n}\)，矛盾。因此 \(\mathfrak{q} = \mathfrak{n}^0\)。

推论 1。令 \(\mathfrak{q}\) 为异于 \(\mathfrak{g}\) 的 \(\mathfrak{g}\) 的子代数集合中的极大元素。则 \(\mathfrak{q}\) 或者是抛物的，或者在 \(\mathfrak{g}\) 中约化。

我们可以假设 g 是某个有限维向量空间 V 的 \(\mathfrak{gl}(V)\) 子代数。令 \(\mathfrak{e}(\mathfrak{q}) \subset \mathfrak{g}\) 为 q 的可分解包络。若 \(\mathfrak{e}(\mathfrak{q}) = \mathfrak{g}\)，则 q 是 g 的理想（第 VII 章第5节第2号命题4），因而是半单的，从而 q 在 g 中约化。假设 \(\mathfrak{e}(\mathfrak{q}) \neq \mathfrak{g}\)。则 \(\mathfrak{e}(\mathfrak{q}) = \mathfrak{q}\)，所以 q 可分解。假设 q 在 g 中不约化。令 n 为 q 的根基中的幂零元素所成的集合。则 \(n \neq 0\)（第 VII 章第5节第3号命题7 (i)）。令 p 为 n 在 g 中的正规化子。则 \(p \supset q\)，且由于 g 是半单的，\(p \neq g\)。于是 p = q。因此 q 是抛物的（定理1）。

推论 2。令 \(\mathfrak{n}\) 为由幂零元素组成的 \(\mathfrak{g}\) 的子代数。存在 \(\mathfrak{q}\) 的一个抛物子代数 \(\mathfrak{g}\)，具有以下性质：

\begin{enumerate}
\def\labelenumi{(\roman{enumi})}
\item
  \(\mathfrak{n}\subset \mathfrak{n}_{\mathfrak{g}}(\mathfrak{q});\)
\item
  \(\mathfrak{n}\) 在 \(\mathfrak{g}\) 中的正规化子包含于 \(\mathfrak{q}\)；
\item
  每个保持 n 不变的 g 的自同构也保持 q 不变。
\end{enumerate}

若 \(\mathfrak{g}\) 可分裂，则 \(\mathfrak{n}\) 包含于 \(\mathfrak{g}\) 的一个 Borel 子代数中。

令 \(\mathfrak{q}_1\) 为 \(\mathfrak{n}\) 在 \(\mathfrak{g}\) 中的正规化子。这是 \(\mathfrak{g}\) 的可分解子代数。令 \(\mathfrak{n}_1 = \mathfrak{n}_{\mathfrak{g}}(\mathfrak{q}_1)\)。归纳地定义 \(\mathfrak{q}_i\) 为 \(\mathfrak{n}_{i-1}\) 在 \(\mathfrak{g}\) 中的正规化子，并令 \(\mathfrak{n}_i\) 等于 \(\mathfrak{n}_{\mathfrak{g}}(\mathfrak{q}_i)\)。序列 \((\mathfrak{n}, \mathfrak{n}_1, \mathfrak{n}_2, \ldots)\) 和 \((\mathfrak{q}_1, \mathfrak{q}_2, \ldots)\) 都是递增的。存在 \(j\) 使 \(\mathfrak{q}_j = \mathfrak{q}_{j+1}\)，换言之，\(\mathfrak{q}_j\) 是 \(\mathfrak{n}_{\mathfrak{g}}(\mathfrak{q}_j)\) 在 \(\mathfrak{g}\) 中的正规化子。因此 \(\mathfrak{q}_j\) 是抛物的（定理1）。我们有 \(\mathfrak{n} \subset \mathfrak{n}_j =\) \(\mathfrak{n}_{\mathfrak{g}}(\mathfrak{q}_{j})\)，且 \(q_{1} \subset q_{j}\)；每个保持 n 不变的 g 的自同构显然都保持 \(n_{1}, n_{2}, \ldots\) 和 \(q_{1}, q_{2}, \ldots\) 不变。若 g 可分裂，则 \(q_{j}\) 包含一个 Borel 子代数 b，从而（§3，第4号，命题13），有 \(\mathfrak{b} \supset \mathfrak{n}_{\mathfrak{g}}(\mathfrak{q}_{j}) \supset \mathfrak{n}\)。

定理 3。假设 k 是代数闭域。令 g 为半单李代数，a 为 g 的可解子代数。则存在包含 a 的 g 的一个 Borel 子代数。

由第 VII 章第5节第2号命题4的推论1 (ii)，可假设 \(\mathfrak{a}\) 是可分解的。存在 g 的一个交换子代数 \(\mathfrak{t}\)，由半单元素组成，使得 \(\mathfrak{g}\)\(\mathfrak{a}\) 是 \(\mathfrak{t}\) 与 \(\mathfrak{n}_{\mathfrak{g}}(\mathfrak{a})\) 的半直积（第 VII 章第5节第3号命题6的推论2）。存在（定理2的推论2）g 的一个抛物子代数 \(\mathfrak{q}\)，使得 \(\mathfrak{g}\)\(\mathfrak{n}_{\mathfrak{g}}(\mathfrak{a}) \subset \mathfrak{n}_{\mathfrak{g}}(\mathfrak{q})\)，且 \(\mathfrak{n}_{\mathfrak{g}}(\mathfrak{a})\) 在 g 中的正规化子包含于 \(\mathfrak{g}\)\(\mathfrak{q}\)；更不用说，\(\mathfrak{a} \subset \mathfrak{q}\)。令 \(\mathfrak{b}\) 为包含于 \(\mathfrak{g}\)\(\mathfrak{q}\) 的 g 的一个 Borel 子代数，令 \(\mathfrak{h}\) 为包含于 \(\mathfrak{g}\)\(\mathfrak{b}\) 的 g 的一个嘉当子代数。于是 \(\mathfrak{h}\) 是 \(\mathfrak{q}\) 的嘉当子代数，所以存在 \(s \in \mathrm{Aut}_e(\mathfrak{q})\)，使 \(s(\mathfrak{t}) \subset \mathfrak{h}\)（第 VII 章第2节第3号命题10和第 VII 章第3节第2号定理1）。我们有 \(s(\mathfrak{n}_{\mathfrak{g}}(\mathfrak{q})) = \mathfrak{n}_{\mathfrak{g}}(\mathfrak{q})\)（第 VII 章第3节第1号注），所以

\[
s (\mathfrak {a}) = s (t) + s \left(\mathfrak {n} _ {\mathfrak {g}} (\mathfrak {a})\right) \subset \mathfrak {h} + s \left(\mathfrak {n} _ {\mathfrak {g}} (\mathfrak {q})\right) = \mathfrak {h} + \mathfrak {n} _ {\mathfrak {g}} (\mathfrak {q}) \subset \mathfrak {b}.
\]

推论。若 \(k\) 是代数闭域，则 \(\mathfrak{g}\) 的每个极大可解子代数都是 Borel 子代数。

